% Source-derived chapter 04: Compactification of the Zastava model
% Original source: intersection-complexes-unramified-l-factors
\chapter{Compactification of the Zastava model}
\label{intersection-complexes-unramified-l-factors-chapter-04}
\source{intersection-complexes-unramified-l-factors}

\begin{remark}[Source section]
\label{intersection-complexes-unramified-l-factors-chapter-04-source}
\source{intersection-complexes-unramified-l-factors}
\source{intersection-complexes-unramified-l-factors}
This chapter reproduces the corresponding section of the retrieved
LaTeX source, including its definitions, statements, examples, and
proofs. It has not yet been translated into Lean.
\notready
\end{remark}

% The source section title was: Compactification of the Zastava model
\label{sect:compactification}
\source{intersection-complexes-unramified-l-factors}

The map $\pi : \Scr Y \to \Scr A$ defined in \eqref{mappi}
is in general not proper, so for example
we cannot apply the decomposition theorem. To rectify this, we introduce
a compactification.

\subsection{Basic properties} \label{sect:basicproperties-comp}
\source{intersection-complexes-unramified-l-factors}

Let $\ol{G/N} = \Spec k[G/N]$ denote the canonical affine closure of the 
quasiaffine variety $G/N$. 
For an arbitrary connected reductive group $G$, Drinfeld's compactification 
$\ol\Bun_B$ is defined\footnote{The definition as a closure is only true in characteristic $0$. In positive characteristic, see \cite[\S 4.1]{ABB}, \cite[\S 7.2]{Sch}.}
 as the closure of $\Bun_B$ inside 
\[ \Maps_\gen(C, G \bs \ol{G/N}/T \supset \pt/B), \] 
the stack parametrizing maps $C \to G\bs \ol{G/N}/T$ that generically land in 
the open substack $G \bs (G/N)/T = \pt/B$.
(When $[G,G]$ is simply connected, 
\cite[Proposition 1.2.3]{BG} show that $\Bun_B$ is dense in 
$\Maps_\gen(C, G\bs \ol{G/N}/T \supset \pt/B)$.) 


Consider the \emph{stack} quotient $X \xt^G \ol{G/N} := (X\times \ol{G/N})/G$, where $G$ acts 
anti-diagonally. Then $X/N = X\times^G G/N$ is an open substack of 
$X \xt^G \ol{G/N}$, so we also have the open substack 
$\pt = X^\circ/B \subset X/B \subset X \xt^G \ol{G/N}/T$. 
Define 
\begin{equation} \label{e:defbarY}
\source{intersection-complexes-unramified-l-factors}
    \barY = (\sM_X \xt_{\Bun_G} \ol \Bun_B)^\circ \subset 
    \Maps_\gen( C, X \xt^G \ol{G/N}/T \supset \pt)  
\end{equation}
where the superscript $^\circ$ denotes the open substack of 
\[ \sM_X \xt_{\Bun_G} \ol \Bun_B \subset \Maps_\gen( C, X \xt^G \ol{G/N}/T \supset 
X^\bullet / B) \] 
parametrizing maps generically landing in $\pt = X^\circ/B$. 
In particular, $\barY$ is an algebraic stack locally of finite type. 
We can identify
$\Scr Y \cong \barY \xt_{\ol\Bun_B} \Bun_B$ as an open substack of $\barY$.
(If $[G,G]$ is simply connected, the containment in \eqref{e:defbarY} is an 
equality.)

There is a natural map from $X\xt^G \ol{G/N}$ to 
\[ (X \times \ol{G/N})/\!\!/G = \Spec k[X \times G]^{G \times N}. \]
Since $k[X \times G]^G = k[X]$, we deduce that $(X \times \ol{G/N})/\!\!/G = X/\!\!/N$.
Therefore, we have a map 
$X \xt^G \ol{G/N} \to X/\!\!/N$ extending the natural map $X/N \to X/\!\!/N$.
Applying $\Maps(C,?/T)$ to the former, we have constructed a map 
\begin{equation}
    \bar\pi : \barY \to \Scr A 
\end{equation}
extending $\pi : \Scr Y \to \Scr A$.
Let $\barY^{\ch\lambda}$ denote the preimage of the subscheme $\Scr A^{\ch\lambda}$. 
The same proof as in Proposition~\ref{prop:factorization} 
shows that $\barY$ has the graded factorization property. 
We will see from Lemma~\ref{lem:barYGr} below that $\barY^{\ch\lambda}$ is 
representable by a scheme of finite type over $k$.

\medskip


First, we show that $\bar\pi$ is indeed a compactification:

\begin{prop}
\source{intersection-complexes-unramified-l-factors}
\notready \label{prop:piproper}
\source{intersection-complexes-unramified-l-factors}
The map $\bar \pi :\barY\to \Scr A$ is proper.
\end{prop}

We proceed as in \S\ref{sss:YGr}. 
Choose $\delta\in \Lambda_X$ lying on the interior of the cone 
dual to $\Cal C_0(X)$.
Then $\delta$ defines a map $X/\!\!/N \to \mbb A^1$, which induces a map 
$\Scr A \to \Sym C$. 
This allows us to consider $\Scr Y$ as a scheme over $\Sym C$.
Proposition~\ref{prop:YGr} gives
a closed embedding $\Scr Y \into \Gr_{B,\Sym C}$ over $\Sym C$. 
We compose this with the natural map $\Gr_{B,\Sym C} \to \Gr_{G,\Sym C}$
to get a map $\Scr Y \to \Gr_{G,\Sym C}$. 
We can extend this to a map 
\begin{equation} \label{e:barYGr}
\source{intersection-complexes-unramified-l-factors}
\barY \to \Gr_{G,\Sym C} 
\end{equation}
using the same idea as in the definition of \eqref{e:YGr}. Namely, 
let $y: C\times S \to X \times^G \ol{G/N} / T$ be an $S$-point of $\barY$
and let $D\subset C\times S$ denote the divisor it maps to. 
In particular, $y$ defines a $G$-bundle $\Scr P_G$ on $C\times S$ 
with a $B$-reduction on $C\times S\sm D$. Since $y(C\times S\sm D) = X^\circ/B =\pt$,
the $G$-bundle in fact admits a trivialization on $C\times S\sm D$. 
The data of $D, \Scr P_G$, and the trivialization defines an $S$-point of $\Gr_{G,\Sym C}$.

\begin{lem}
\source{intersection-complexes-unramified-l-factors}
\notready \label{lem:barYGr}
\source{intersection-complexes-unramified-l-factors}
The map $\barY \to \Gr_{G,\Sym C} \underset{\Sym C}\times \Scr A$
is a closed embedding.
\end{lem}

Since $\Gr_{G,\Sym C}$ is ind-proper over $\Sym C$ (cf.~\cite[Remark 3.1.4]{Xinwen}), 
Proposition~\ref{prop:piproper} follows from Lemma~\ref{lem:barYGr}.
We also deduce from the lemma that $\barY$ is representable by a scheme, and 
each $\barY^{\ch\lambda}$ is of finite type.


\begin{proof}
The discussion above really defines a map 
\begin{equation} \label{e:biggeremb}
\source{intersection-complexes-unramified-l-factors}
    \Maps_\gen(C, X \xt^G \ol{G/N}/T \supset \pt) \to \Gr_{G,\Sym C} \xt_{\Sym C} \sA, 
\end{equation}
and $\barY \into \Maps_\gen(C, X\xt^G \ol{G/N}/T \supset \pt)$ is a closed
embedding. Thus, it suffices to prove that \eqref{e:biggeremb} is a closed
embedding.  
Fix a test scheme $S$. 
Let $\Scr P^0_G,\Scr P^0_B,\Scr P^0_T$ denote the respective trivial bundles on $C\times S$. 
An $S$-point of $\Gr_{G,\Sym C} \times_{\Sym C} \Scr A$
consists of the data 
\[ (\Scr P_G,\Scr P_T,D,\tau,\alpha)  \]
with $\Scr P_G\in \Bun_G(S),\, \Scr P_T \in \Bun_T(S),\, D\in \Sym C(S)$,
a trivialization $\tau : \Scr P^0_G|_{C\times S\sm D}\cong \Scr P_G|_{C\times S\sm D}$, 
and a section
$\alpha: C\times S \to (X/\!\!/N)\times^T \Scr P_T$ such that 
$(\Scr P_T,\alpha) \in \Scr A(S)$ maps to $D$. 
In particular, this means that $\alpha$ induces a trivialization 
$\Scr P_T^0|_{C\times S\sm D} \cong \Scr P_T|_{C\times S\sm D}$. 
Using the identification $X^\circ \cong B$, we get a section
\[ \sigma_0 : C\times S\sm D \to \Scr P^0_B \cong X^\circ \xt^B \Scr P_B^0 
\into X \xt^B \Scr P_B^0= X \xt^G \Scr P_G^0. \]
The composition $\sigma := \tau \circ \sigma_0$ then defines a section
$C\times S \sm D \to X\xt^G \Scr P_G$.
On the other hand, the trivial $B$-bundle also corresponds to a section
$\kappa_0 : C\times S \sm D \to \Scr P_G^0 \xt^G (G/N)$.
Composing $\kappa_0$ with the trivialization $\alpha:\Scr P_T^0|_{C\times S\sm D} \cong \Scr P_T|_{C\times S\sm D}$, we get a section
\[ \kappa : C\times S \sm D \to \Scr P_G \xt^G \ol{G/N} \xt^T \Scr P_T. \]
The datum $(\Scr P_G, \Scr P_T, \sigma, \kappa)$ defines an $S$-point 
of $\Maps(C, X\xt^G \ol{G/N}/T)$ if and only if $\sigma,\kappa$ both extend to regular
maps on $C\times S$. 
Therefore, the fiber of our chosen $S$-point over 
the map \eqref{e:biggeremb}
parametrizes maps $S' \to S$
such that the base change of $\sigma,\kappa$ to $S'$ both extend to 
$C\times S'$. By Lemma~\ref{lem:extendsection}, 
this fiber is represented by a closed subscheme of $S$ (here the key point 
is that both $X$ and $\ol{G/N}$ are affine). 
\end{proof}

\begin{rem}
\source{intersection-complexes-unramified-l-factors}
\notready
We need the extra factor of $\Scr A$ in Lemma~\ref{lem:barYGr} which was
not present in Proposition~\ref{prop:YGr} because the map $\Gr_B \to \Gr_G$
is a bijection on $k$-points but far from an isomorphism of ind-schemes. 
Since $\Scr A$ embeds into $\Gr_{T,\Sym C}$, the lemma is really
embedding $\barY$ into $\Gr_{G,\Sym C}\times_{\Sym C} \Gr_{T,\Sym C}$. 
The ind-scheme $\Gr_T$ is highly non-reduced, while $(\Gr_T)_\red$ is a disjoint
union of points.
\end{rem}

\begin{eg}
\source{intersection-complexes-unramified-l-factors}
\notready \label{eg:Heckecompact}
\source{intersection-complexes-unramified-l-factors}
Let $X = \mbb G_m \bs \GL_2$ as in Example~\ref{eg:YHecke}. 
Then $\barY = \Sym C \xt \Sym C$ and $\barY \to \sA$ is the identity morphism. 
\end{eg}

\subsection{Stratification} 
\label{sect:strat-compactified}
\source{intersection-complexes-unramified-l-factors}
We can uniquely write any $\ch\nu \in \ch\Lambda^\pos_G$ 
as a sum $\ch\nu = \sum_{\alpha\in \Delta_G} n_\alpha \ch\alpha$ where $\Delta_G$ is the set of 
simple coroots and $n_\alpha$ are positive integers. 
Let $C_{\ch\nu} := \prod_{\Delta_G} C^{(n_\alpha)}$ denote the corresponding partially
symmetrized power of $C$. 
Recall that Drinfeld's 
compactification $\ol\Bun_B$ of $\Bun_B$ has a 
stratification by \emph{defect}, where the strata are given by 
locally closed embeddings
\[ \mf i_{\ch\nu}: C_{\ch\nu} \times \Bun^{\ch\mu+\ch\nu}_B \into \ol\Bun^{\ch\mu}_B, \]
for $\ch\nu\in \ch\Lambda^\pos_G,\, \ch\mu \in \ch\Lambda_G$, cf.~\cite[\S 1.5, p.~7]{BFGM}. 
Define the substack ${_{\ch\nu}}\ol\Bun_B^{\ch\mu}$ to be the image 
of the corresponding embedding. 
We obtain an open substack $_{\le \ch\nu}\ol\Bun_B^{\ch\mu} \subset \ol\Bun^{\ch\mu}_B$ by taking the union of the strata $_{\ch\nu'}\ol\Bun_B^{\ch\mu}$ 
for all $\ch\nu' \le \ch\nu$. 

Since $\barY^{\ch\lambda}$ maps to $\ol\Bun_B^{-\ch\lambda}$ for $\ch\lambda \in \mathfrak c_X$, by base change we have locally
closed subschemes 
\[ {}_{\ch\nu} \barY^{\ch\lambda} := \barY^{\ch\lambda} \xt_{\ol\Bun_B^{-\ch\lambda}} {}_{\ch\nu} \ol\Bun_B^{-\ch\lambda} \into \barY^{\ch\lambda} \]
and open subschemes ${}_{\le \ch\nu} \barY^{\ch\lambda} \into \barY^{\ch\lambda}$ defined analogously.
Observe that the identification ${}_{\ch\nu} \ol\Bun_B^{-\ch\lambda} \cong 
C_{\ch\nu} \times \Bun_B^{\ch\nu-\ch\lambda}$ induces an isomorphism
\begin{equation} \label{e:barYstrata1}
\source{intersection-complexes-unramified-l-factors}
    {}_{\ch\nu} \barY^{\ch\lambda}  \cong C_{\ch\nu} \times \Scr Y^{\ch\lambda-\ch\nu}. 
\end{equation}
On the other hand, $\barY^{\ch\lambda}$ also maps to $\Scr M_X$, so 
we get a locally closed subscheme 
\[ {}_{\ch\nu}\barY^{\ch\lambda,\ch\Theta} := {}_{\ch\nu}\barY^{\ch\lambda} \xt_{\Scr M_X} \Scr M_X^{\ch\Theta} \into {}_{\ch\nu}\barY^{\ch\lambda} \into \barY^{\ch\lambda} \] 
for $\ch\Theta$ any partition in $\mathfrak c_X^-$ 
(by Lemma~\ref{lem:globstrata}). 
We deduce from \eqref{e:barYstrata1} that there is an isomorphism
\[ {}_{\ch\nu} \barY^{\ch\lambda,\ch\Theta} \cong C_{\ch\nu} \times \Scr Y^{\ch\lambda-\ch\nu, \ch\Theta}. \] 
In particular, Proposition~\ref{prop:localstrata} implies that 
${}_{\ch\nu} \barY^{\ch\lambda,\ch\Theta}$ is smooth. 
In summary: 

\begin{prop}
\source{intersection-complexes-unramified-l-factors}
\notready  \label{prop:barYstrata}
\source{intersection-complexes-unramified-l-factors}
The collection of locally closed subschemes 
${}_{\ch\nu} \barY^{\ch\lambda,\ch\Theta}$, ranging over all $\ch\nu\in \ch\Lambda^\pos_G$ and partitions $\ch\Theta \in \Sym^\infty(\mathfrak c_X^-\sm 0)$, 
forms a smooth stratification of $\barY^{\ch\lambda}$. 
\end{prop}

Note that for fixed $\ch\lambda$, many of these strata may be empty.


\subsubsection{Changing the curve}

In this subsection we let $C$ be a smooth curve which is not necessarily proper 
and we define 
\[ \sA(C) = \Hom( \mf c_X^\vee, \Sym C), \quad 
\sY(C)=\Maps_\gen(C, X/B\supset \pt) \]
and $\barY(C)$ to be the closure
of $\sY(C)$ in $\Maps_\gen(C, X \xt^G \ol{G/N}/T \supset \pt)$ to emphasize
the curve being used. 
Here $\Hom$ denotes homomorphisms of monoid objects in the category of schemes.
Similarly we have $\sA^{\ch\lambda}(C), \sY^{\ch\lambda}(C),\barY^{\ch\lambda}(C)$. 
The local nature of $\barY(C)$ 
(in particular Lemma~\ref{lem:barYGr}) ensures that $\barY^{\ch\lambda}(C)$ is still 
a finite type $k$-scheme. 

\smallskip

Let $p: \wt C \to C$ be an \'etale map of smooth curves. 
Let $(\Sym \wt C)_\disj \subset \Sym \wt C$
be the open subset that consists of divisors $\wt D$ on $\wt C$ such that
any fiber of $p$ contains at most one point of the support of $\wt D$.
Let $\sA(\wt C)_\disj$ denote the open subset of 
$\Hom(\mf c_X^\vee, \Sym \wt C)$ consisting of homomorphisms landing in $(\Sym \wt C)_\disj$.
Then pushforward of divisors defines a map 
$\sA^{\ch\lambda}(\wt C) \to \sA^{\ch\lambda}(C)$.

\begin{prop}[{\cite[Proposition 2.19]{BFG}}]
\source{intersection-complexes-unramified-l-factors}
\notready \label{prop:changecurve}
\source{intersection-complexes-unramified-l-factors}
For an \'etale map $\wt C \to C$ we have a canonical isomorphism 
\[ \barY^{\ch\lambda}(C) \xt_{\sA^{\ch\lambda}(C)} \sA^{\ch\lambda}(\wt C)_\disj 
\cong \barY^{\ch\lambda}(\wt C) \xt_{\sA^{\ch\lambda}(\wt C)} \sA^{\ch\lambda}(\wt C)_\disj 
\] 
which preserves the fine stratification. 
\end{prop}
Note that setting $\wt C = C\sqcup C$ recovers the graded factorization property. 

\begin{proof}
Choose $\delta \in \mf c_X^\vee$ as in \S\ref{sss:YGr}. 
For $(y,\tilde a) \in \barY^{\ch\lambda}(C) \xt_{\sA^{\ch\lambda}(C)} \sA^{\ch\lambda}(\wt C)_\disj$, let $\wt D$ (resp.~$D$) denote the divisor corresponding to $\delta$ paired with 
$\tilde a$ (resp.~$\pi(y)\in \sA^{\ch\lambda}(C)$). 
Then $p_*\wt D = D$ and we deduce that there is an isomorphism 
$\wh C'_D \cong \wh{\wt C}{}'_{\wt D}$.
Lemma~\ref{lem:B-L}, applied to
the affine $G\xt T$-scheme $X \xt \ol{G/N}$, implies that the point 
$y\in \barY^{\ch\lambda}(C)$ is 
equivalent to its restriction $y|_{\wh C'_D}$. 
Applying the same lemma again shows that $y|_{\wh C'_D}$ is equivalent to 
a point $\tilde y \in \barY^{\ch\lambda}(\wt C)$ such that $\tilde y(\wt C \sm \wt D)=\pt$.
This defines mutually inverse maps in both directions and the compatibility 
with strata is clear.
\end{proof}

Since the diagonal $\delta^{\ch\lambda} : \wt C \into \sA^{\ch\lambda}(\wt C)$ 
sending $\tilde v \mapsto \ch\lambda\cdot \tilde v$ is contained in $\sA^{\ch\lambda}(\wt C)_\disj$, we deduce from the graded factorization property
and Proposition~\ref{prop:changecurve} that $\barY^{\ch\lambda}(C)$ 
is \'etale-locally isomorphic to $\barY^{\ch\lambda}(\bbA^1)$ for any smooth curve $C$.


\subsection{The central fiber} \label{def:centralfiber}
\source{intersection-complexes-unramified-l-factors}
Let $\ch\lambda \in \mathfrak c_X$. There is a 
diagonal map $\delta^{\ch\lambda}:C \to \Scr A^{\ch\lambda}$
sending $v \mapsto \ch\lambda \cdot v$. For a fixed point
$v\in \abs C$, let us consider 
$\delta^{\ch\lambda}_v : v \to \Scr A^{\ch\lambda}$. 

\medskip

Define the central fiber $\msf Y^{\ch\lambda}$ of $\Scr Y^{\ch\lambda}$ 
to be the
preimage of $\delta^{\ch\lambda}_v$ under 
$\pi : \Scr Y^{\ch\lambda} \to \Scr A^{\ch\lambda}$. 
If we take central fibers of the map \eqref{e:YGr}, Proposition~\ref{prop:YGr} implies that we have a closed embedding 
$\msf Y^{\ch\lambda} \into \Gr_B$.
In fact, $\ch\lambda \cdot v$ can be considered as a point in 
$\Gr_T$. If we let $\msf S^{\ch\lambda}$ denote the preimage of this point
under the projection $\Gr_B \to \Gr_T$, then we get a closed embedding 
\begin{equation} \label{e:YembedGr}
\source{intersection-complexes-unramified-l-factors}
    \msf Y^{\ch\lambda} \into \msf S^{\ch\lambda}.
\end{equation}

Note that $\msf S^{\ch\lambda}$ identifies with 
the $\msf L N$-orbit of $t^{\ch\lambda}$ in $\Gr_G$. 
The orbits $\msf S^{\ch\lambda}$ are commonly known
as the semi-infinite orbits of $\Gr_G$, and their geometric properties
were extensively studied by Mirkovi\'c--Vilonen in \cite{MV}. 
Let $\olsf S^{\ch\lambda}$ denote the scheme-theoretic closure 
of $\msf S^{\ch\lambda}$ in $\Gr_G$.

We analogously define $\olsf Y^{\ch\lambda}$ as the central fiber\footnote{The open subscheme $\msf Y^{\ch\lambda}$ does not need to be dense in $\olsf Y^{\ch\lambda}$.}
of $\bar\pi : \barY^{\ch\lambda} \to \sA^{\ch\lambda}$, and 
we deduce from Lemma~\ref{lem:barYGr} that there is a 
closed embedding $\olsf Y^{\ch\lambda} \into \Gr_G$. 
From the moduli-theoretic description of $\olsf S^{\ch\lambda}$ 
(see \cite[proof of Proposition 5.3.6]{Xinwen}), we see that
this factors through a closed embedding 
\begin{equation} \label{e:olYembedGr}
\source{intersection-complexes-unramified-l-factors}
    \olsf Y^{\ch\lambda} \into \olsf S^{\ch\lambda}.
\end{equation}

\subsubsection{Local description of $\msf Y^{\ch\lambda}$} 
Note that the central fiber $\msf Y^{\ch\lambda}$
intersects the stratum $\Scr Y^{\ch\lambda,\ch\Theta}$ 
only if $\ch\Theta =[\ch\theta]$ is the singleton partition corresponding to 
a single $\ch\theta \in \mathfrak c_X^-$. 
In this case, let $\msf Y^{\ch\lambda,\ch\theta}$ denote
the intersection $\msf Y^{\ch\lambda} \cap \Scr Y^{\ch\lambda,\ch\theta} = \msf Y^{\ch\lambda}
\xt_{\sM_X} \sM_X^{\ch\theta}$.
Also let $\olsf Y^{\ch\lambda,\ch\theta} = \olsf Y^{\ch\lambda} \xt_{\sM_X} \sM_X^{\ch\theta}$.

The $\msf L G$-action on the base point $x_0 \in \msf L X(k)$
defines a map $\msf L G \to \msf L X$, which induces a map 
$\Gr_G \to \msf L X/ \msf L^+G$.  

\begin{lem}
\source{intersection-complexes-unramified-l-factors}
\notready \label{lem:Y=ScapGr}
\source{intersection-complexes-unramified-l-factors}
There are natural isomorphisms
\begin{align} 
    \msf Y^{\ch\lambda} & 
        \cong \msf S^{\ch\lambda} \xt_{\msf L X/\msf L^+G} \msf L^+ X/\msf L^+G 
\label{e:YScapGr} \\
\source{intersection-complexes-unramified-l-factors}
    \olsf Y^{\ch\lambda} & 
        \cong \olsf S^{\ch\lambda} \xt_{\msf L X/\msf L^+G} \msf L^+ X / \msf L^+G 
\label{e:olYScapGr}
\source{intersection-complexes-unramified-l-factors}
\end{align}
which induce 
$\displaystyle (\msf Y^{\ch\lambda,\ch\theta})_\red \cong (\msf S^{\ch\lambda} \xt_{\msf L X/\msf L^+G} \msf L^{\ch\theta}X/\msf L^+G)_\red$ 
and 
$\displaystyle (\ol{\msf Y}^{\ch\lambda,\ch\theta})_\red \cong (\ol{\msf S}^{\ch\lambda}
 \xt_{\msf L X/\msf L^+G} \msf L^{\ch\theta}X/\msf L^+G)_\red$.
\end{lem}


\begin{proof}
Consider the composition $\Gr_B^{\ch\lambda} \to  \Gr_G \to \msf L X/\msf L^+G$. 
It follows from the definitions that we have an embedding
\[  \msf Y^{\ch\lambda} \into \Gr^{\ch\lambda}_B \xt_{\msf L X/\msf L^+G} \msf L^+X/\msf L^+G
\]
The map in the reverse direction is defined using Beauville--Laszlo's theorem: for a $k$-algebra $R$, an $R$-point of $\Gr^{\ch\lambda}_B$ consists
of a $B$-bundle $\hat{\Scr P}_B$ on $\Spec R\tbrac t$ and a section 
$\hat\sigma_0 : \Spec R\lbrac t \to \hat{\Scr P}_B$. Using the
identification $B \cong X^\circ$, we can identify $\hat\sigma_0$
with a section $\hat\sigma: \Spec R\lbrac t \to X^\circ \times^B \hat{\Scr P}_B$. 
If the image of $(\hat{\Scr P}_B, \hat \sigma)$ in $\msf L X/\msf L^+G$ belongs to $\msf L^+X/\msf L^+G$, then 
$\hat\sigma$ extends to a section $\Spec R\tbrac t \to X\times^B \hat{\Scr P}_B$. 
By Lemma~\ref{lem:B-L}, the pair $(\hat{\Scr P}_B, \hat \sigma)$ is equivalent to an $R$-point $y\in \Scr Y$. 
By construction, $y\in \msf Y^{\ch\lambda}$.
The two maps are mutually inverse, so we get \eqref{e:YScapGr}.

The same argument proves \eqref{e:olYScapGr},
and the other identifications are by definition.
\end{proof}

We included the subscript $_\red$ above for clarity, but from now on we 
consider only the underlying reduced structure on all central fibers 
and omit the subscript since the \'etale site is insensitive
to reduced structures.

\begin{eg}
\source{intersection-complexes-unramified-l-factors}
\notready \label{eg:Yfiber}
\source{intersection-complexes-unramified-l-factors}
Resume the setting of Example~\ref{eg:YHecke}.  
Then $\msf Y^{n_1 \ch\nu_{D^+} + n_2 \ch\nu_{D^-}}$ is empty if both $n_1$ and $n_2$ are nonzero.
Otherwise it consists of a single point. 
 If we use the embedding $\mbb G_m \into \GL_2$
via $a\mapsto \left( \begin{smallmatrix} 1 & 0 \\ 1-a& a \end{smallmatrix}\right)$
and fix the base point $x_0=1$, then $\msf Y^{\ch\nu_{D^+}}$ 
corresponds to the point 
$\left( \begin{smallmatrix} 1 & 1 \\ 0 & 1 \end{smallmatrix} \right) t^{\ch\vareps_1} \in \Gr_B$ 
while $\msf Y^{\ch\nu_{D^-}}$ corresponds to $t^{-\ch\vareps_2} \in \Gr_B$.
\end{eg}


\subsubsection{} \label{sect:YttimesC}
\source{intersection-complexes-unramified-l-factors}
Lemma~\ref{lem:Y=ScapGr} implies that 
the central fibers $\msf Y^{\ch\lambda}, \olsf Y^{\ch\lambda}$ 
can be defined purely locally (in particular, independently of 
the point $v\in \abs C$).
We also deduce that $\sY^{\ch\lambda} \xt_{\sA^{\ch\lambda},\delta^{\ch\lambda}} C 
\cong \msf Y^{\ch\lambda} \ttimes C$ and 
$\barY^{\ch\lambda} \xt_{\sA^{\ch\lambda},\delta^{\ch\lambda}} C \cong \olsf Y^{\ch\lambda} \ttimes C$,
where $? \ttimes C := {?}\times^{\Aut k\tbrac t} C^\wedge $
and $C^\wedge \to C$ is the $\Aut k\tbrac t$-torsor classifying $v\in C$
together with an isomorphism $\mf o_v \cong k\tbrac t$. 


\subsection{Dimension of central fibers}

We will now discuss an argument that is critical when estimating dimensions of Zastava spaces and their central fibers. This argument appeared in the proof 
of \cite[Theorem 3.2]{MV}. The second author thanks M.~Finkelberg for
explaining this proof to him.

Semi-infinite orbits have a very simple geometric structure, summarized by the following:
\begin{prop}[{\cite[Proposition 3.1]{MV}, \cite[Proposition 5.3.6, Corollary 5.3.8]{Xinwen}}]
\source{intersection-complexes-unramified-l-factors}
\notready  \label{prop:Shyperplane}
\source{intersection-complexes-unramified-l-factors}
\ 

\begin{enumerate}
\item We have a stratification
$\olsf S^{\ch\lambda} = \bigcup_{\ch\lambda' \le \ch\lambda} \msf S^{\ch\lambda'}$.

\item
Inside $\olsf S^{\ch\lambda}$, the boundary of $\msf S^{\ch\lambda}$
is given by a hyperplane section under an embedding of $\Gr_G$ in projective space. 
\end{enumerate}
\end{prop}



In particular, if a projective subvariety meets the semi-infinite orbit $\msf S^{\ch\lambda}$, it also meets its boundary $\bigcup_{\lambda' < \ch\lambda}  \msf S^{\ch\lambda'}$. We will use this simple fact several times, in order to estimate the dimensions of central fibers.

\medskip

For $\ch\lambda \in \mf c_X$, consider
the central fiber $\olsf Y^{\ch\lambda} \subset \olsf S^{\ch\lambda} \subset \Gr_G$. 
Observe that the defect stratification of $\barY^{\ch\lambda}$ (\S \ref{sect:strat-compactified}) gives a stratification
of $\olsf Y^{\ch\lambda} = \bigcup_{\ch\lambda' \le \ch\lambda} \msf Y^{\ch\lambda'}$,
where $\msf Y^{\ch\lambda'} = \olsf Y^{\ch\lambda}\cap ({_{\ch\lambda-\ch\lambda'}} \barY^{\ch\lambda})$. 
This is compatible with the stratification of $\olsf S^{\ch\lambda} = \bigcup_{\ch\lambda'\le\ch\lambda}
\msf S^{\ch\lambda'}$ under the closed embedding $\olsf Y^{\ch\lambda}\subset \olsf S^{\ch\lambda}$.
From Lemma~\ref{lem:Y=ScapGr} we have $\olsf Y^{\ch\lambda} \cap \msf S^{\ch\lambda'} = \msf Y^{\ch\lambda'}$
for $\ch\lambda' \le \ch\lambda$.


\begin{prop}
\source{intersection-complexes-unramified-l-factors}
\notready
\label{prop:intersections}
\source{intersection-complexes-unramified-l-factors}
 Given an irreducible component $\olsf Y \subset \olsf S^{\ch\lambda}$ of the central fiber $\olsf Y^{\ch\lambda}$ of $\barY^{\ch\lambda}$, there is a $\ch\lambda'\le \ch\lambda$ with $\olsf Y \cap \msf S^{\ch\lambda'}$ nonempty of dimension zero, and $d:=\dim \olsf Y \le \brac{ \rho_G, \ch\lambda-\ch\lambda'}$. 

Moreover, if $d= \brac{ \rho_G, \ch\lambda-\ch\lambda'}$, then there exists a sequence of simple roots $\alpha_1, \dots,\alpha_d$ (possibly with repetitions) 
and subvarieties $\msf Y_j \subset \msf Y^{\ch\lambda-\ch\alpha_1-\dotsb-\ch\alpha_j}$ for $j=0,\dotsc,d$  such that 
\begin{itemize}
\item $\ol{\msf Y_0} = \olsf Y$,
\item $\msf Y_j$ is an irreducible component of $\ol{\msf Y_{j-1}} \cap \msf S^{\ch\lambda - \ch\alpha_1 - \cdots -\ch\alpha_j}$ of dimension $d-j$ for $j=1,\dotsc, d$.
\end{itemize}
\end{prop}

\begin{proof}
By the stratification $\olsf S^{\ch\lambda} = \bigcup_{\ch\lambda'\le\ch\lambda} \msf S^{\ch\lambda}$,
there must exist $\ch\lambda_0 \le \ch\lambda$ such that 
$\msf Y_0 := \olsf Y \cap \msf S^{\ch\lambda_0}$ is dense in $\olsf Y$. 
Then $\olsf Y \subset \olsf S^{\ch\lambda_0}$ so we may assume $\ch\lambda=\ch\lambda_0$. 
Write $\partial \olsf S^{\ch\lambda} = \olsf S^{\ch\lambda} \sm \msf S^{\ch\lambda}$ for the hyperplane of Proposition~\ref{prop:Shyperplane}(ii). 
Since $\olsf Y$ is a projective subvariety of $\olsf S^{\ch\lambda}$, it must
 meet $\partial \olsf S^{\ch\lambda}$. 
Hence there exists $\ch\lambda_1 < \ch\lambda$ such that
$\dim(\olsf Y \cap \msf S^{\ch\lambda_1}) = d-1$. 
Let $\msf Y_1$ be any irreducible component of $\olsf Y\cap \msf S^{\ch\lambda_1}$ and 
$\ol{\msf Y_1}$ its closure in $\olsf S^{\ch\lambda_1}$. 
Continuing in this fashion we produce a sequence of coweights $\ch\lambda_0=\ch\lambda, \ch\lambda_1,\dotsc,\ch\lambda_d$ and irreducible components $\msf Y_j \subset \ol{\msf Y_{j-1}}\cap \msf S^{\ch\lambda_j}$ for $j=1,\dotsc,d$, such that $\ch\lambda_j < \ch\lambda_{j-1}$ and 
$\dim \msf Y_j = d-j$. 
Since $\ch\lambda_j < \ch\lambda_{j-1}$ implies $\brac{\rho_G,\ch\lambda_{j-1}-\ch\lambda_j}\ge 1$, 
we get $\brac{\rho_G,\ch\lambda-\ch\lambda_d} \ge d$ and $\ch\lambda'=\ch\lambda_d$ is the claimed
coweight in the proposition statement. 

In order for the last inequality to be an equality, we must have $\ch\lambda=\ch\lambda_0$
and $\ch\lambda_j = \ch\lambda_{j-1} - \ch\alpha_j$ for some simple coroot $\ch\alpha_j,\, j=1,\dotsc,d$. 
\end{proof}

 


\subsection{Comparison of IC complexes} \label{sect:compareIC}
\source{intersection-complexes-unramified-l-factors}
In order to describe the restriction of $\IC_{\barY}$ to strata 
we will need to introduce several (graded) factorization algebras on the
collection of $C_{\ch\nu},\, \ch\nu\in\ch\Lambda^\pos_G$. 
We only state their definitions; we refer the reader to \cite{BG2, G:whatacts} for
further context.

\smallskip

Let $\ch{\mf n}_C = \ch{\mf n} \ot_{\ol\bbQ_\ell} {\ol\bbQ_\ell}_C$ be the constant sheaf of 
Lie algebras over $C$. 
Let $\mf U(\ch{\mf n}_C)^{\ch\nu}$ denote the following sheaf on $C_{\ch\nu}$
for $\ch\nu \in \ch\Lambda^\pos_G$: its $*$-stalk at a point $\sum \ch\nu_i \cdot v_i$
with $v_i \in \abs C$ distinct is the tensor product $\bigotimes_i U(\ch{\mf n})^{\ch\nu_i}$
where the superscript $\ch\nu_i$ refers to the corresponding weight space of 
the universal enveloping algebra $U(\ch{\mf n})$. 
These stalks glue to a sheaf by means of the co-multiplication map on $U(\ch{\mf n})$. 
Let $\mf U^\vee(\ch{\mf n}_C)^{-\ch\nu} = \mbb D (\mf U(\ch{\mf n}_C)^{\ch\nu}) \in \rmD^b_c(C_{\ch\nu})$ denote the Verdier dual. 

We will also need the Chevalley--Cousin complex $\Upsilon(\ch{\mf n}_C)$ 
of $\ch{\mf n}_C$. 
Consider the (homological) Chevalley complex $C_\bullet(\ch{\mf n}_C)$
as a sheaf of co-commutative DG co-algebras on $C$, endowed with a grading by elements of
$\ch\Lambda^\pos_G$. 
By the general procedure of \cite[\S 3.4]{BD:chiral}, a sheaf of $\ch\Lambda^\pos_G$-graded
co-commutative DG co-algebras on $C$ is equivalent to a commutative factorization algebra 
on $\bigsqcup C_{\ch\nu}$. We let $\Upsilon(\ch{\mf n}_C)^{\ch\nu}$ denote
the corresponding complex on $C_{\ch\nu}$, which can be explicitly constructed
via a Cousin complex. 
Its $*$-stalk at $\sum \ch\nu_i \cdot v_i$ with $v_i$'s disinct is 
the tensor product $\bigotimes_i C_\bullet(\ch{\mf n})^{\ch\nu_i}$. 
We remark that $\Upsilon(\ch{\mf n}_C)^{\ch\nu}$ is actually
a perverse sheaf, and $\Upsilon(\ch{\mf n}_C)$ and $\mf U^\vee(\ch{\mf n}_C)$ are
related by a certain Koszul duality.


\subsubsection{} Let 
\[ \mf i_{\barY,\ch\nu} : C_{\ch\nu} \times \Scr Y^{\ch\lambda-\ch\nu}
\into \barY^{\ch\lambda} \]
denote the locally closed embedding corresponding to \eqref{e:barYstrata1}. 

\begin{prop}
\source{intersection-complexes-unramified-l-factors}
\notready \label{prop:ICbarY}
\source{intersection-complexes-unramified-l-factors}
For any $\ch\lambda\in \mathfrak c_X,\, \ch\nu\in \ch\Lambda^\pos_G$, 
there is an equality 
\[ 
[\mf i^*_{\barY,\ch\nu}(\IC_{\barY^{\ch\lambda}}) ]
= [\mfU^\vee(\ch\mfn_C)^{-\ch\nu} \bt \IC_{\Scr Y^{\ch\lambda-\ch\nu}}] 
\]
in the Grothendieck group of perverse sheaves on $C_{\ch\nu} \xt \sY^{\ch\lambda-\ch\nu}$.
\end{prop}

We will prove the proposition in the course of this subsection. It is based on the following result of {\cite[Proposition 4.4]{BG2}, \cite[Theorem 1.12]{BFGM}}. 
\begin{thm}
\source{intersection-complexes-unramified-l-factors}
\notready
\label{thm:BFGM}
\source{intersection-complexes-unramified-l-factors}
There exists a canonical isomorphism 
\[ \mf i_{\ch\nu}^* (\IC_{\ol\Bun_B^{\ch\mu}}) \cong 
\mfU^\vee(\ch{\mfn}_C)^{-\ch\nu} \bt \IC_{\Bun_B^{\ch\mu+\ch\nu}} 
\]
for any $\ch\nu\in \ch\Lambda^\pos_G,\, \ch\mu\in \ch\Lambda_G$. 
\end{thm}



Given a map $f : Y \to S$ between finite type algebraic $k$-stacks, 
we use the notion of a complex on $Y$ which is 
\emph{universally locally acyclic} (ULA) with 
respect to $f$, as in \cite[D\'efinition 2.12]{SGA4h}.

The general lemma we will use is the following:

\begin{lem}[{\cite[Lemma 7.1.3]{BG}}]
\source{intersection-complexes-unramified-l-factors}
\notready \label{lem:ULAbox} 
\source{intersection-complexes-unramified-l-factors}
Consider a Cartesian diagram of finite type algebraic stacks
\[
\begin{tikzcd} 
Y' \ar[r, "{f'}"] \ar[d, "{g'}", swap] & S' \ar[d, "g"] \\ 
Y \ar[r, "f"] & S
\end{tikzcd}
\]
where $S$ is smooth.
Let $j : Y_0 \into Y$ be an open dense substack such that
the map $f \circ j : Y_0 \to S$ is smooth. 
In addition, assume that the complexes $\IC_Y$ and $j_!(\IC_{Y_0})$
are ULA with respect to the map $f$. 

Denote the closure\footnote{In general it is possible for $Y'$ to have more irreducible
components than $\ol{Y_0\xt_{S} S'}$.} 
of $Y_0 \xt_S S'$ in $Y'$ by $\ol{Y_0 \xt_S S'}$. 
Then there is a natural isomorphism 
\[ 
\IC_{\ol{Y_0 \xt_S S'}} \cong \IC_{S'}\bt_S \IC_{Y}  := 
f'^*(\IC_{S'})\ot g'^*(\IC_Y) [-\dim S],
\]
where the left hand side is implicitly extended by zero to $Y'$.
\end{lem}
For a proof of Lemma~\ref{lem:ULAbox}, see \cite{Wang:ULA}.

\medskip

We would like to apply this lemma with $Y = \ol\Bun_B, \, Y_0 = \Bun_B,\,
S = \Bun_G$, and $S' = \sM_X$. Unfortunately the stack $\ol\Bun_B$ is
``too big'' for all of $\IC_{\ol\Bun_B}$ to be ULA over $\Bun_G$. 
However, we do get the ULA property if we restrict to open substacks
where the defect is not ``too big'':
Let $\mf j : \Bun_B \into \ol\Bun_B$ denote the open embedding.

\begin{prop}[{\cite{Campbell, Campbell2}}]
\source{intersection-complexes-unramified-l-factors}
\notready
\label{prop:ULA}
\source{intersection-complexes-unramified-l-factors}
Fix $\ch\nu \in \ch\Lambda^\pos_G$. Then for any $\ch\mu \in \ch\Lambda_G$
large enough:
\begin{enumerate}
\item The complexes $\IC_{_{\le \ch\nu}\ol\Bun^{-\ch\mu}_B}$ and 
$\mf j_! \IC_{\Bun_B}|_{_{\le \ch\nu}\ol\Bun^{-\ch\mu}_B}$ are 
ULA over $\Bun_G$.
\item The fiber of $\Bun^{-\ch\mu}_B \to \Bun_G$ is dense in the
fiber of $_{\le\ch\nu}\ol\Bun^{-\ch\mu}_B \to \Bun_G$ over any $k$-point of $\Bun_G$.
\end{enumerate}
\end{prop}

Here ``large enough'' is in the same sense as in Lemma~\ref{lem:BunBtoGsm}, i.e.,
deep enough in the dominant chamber $\ch\Lambda^+_G$.

\begin{proof}
The ULA property for $\IC_{_{\le \ch\nu}\ol\Bun^{-\ch\mu}_B}$ 
is \cite[Corollary 4.1.1.1]{Campbell}. 
The ULA property for $\mf j_! \IC_{\Bun_B}|_{_{\le \ch\nu}\ol\Bun^{-\ch\mu}_B}$ can be deduced from the former, cf.~\cite[\S 4.3]{Campbell2}. 
We review the salient features of the proof in \cite{Campbell} to deduce (ii). 

The key object introduced in \cite{Campbell} is Kontsevich's compactification
$\ol\Bun_B^K$ of $\Bun_B$, which is a resolution of singularities $\ol\Bun_B^K \to \ol\Bun_B$
of Drinfeld's compactification. For an affine test scheme $S$, 
an $S$-point of $\ol\Bun_B^K$ is a commutative square
\[ 
\begin{tikzcd}
\wt{\Scr C} \ar[d] \ar[r] & \pt/B \ar[d] \\ C\xt S \ar[r, "\sP_G"] & \pt/G 
\end{tikzcd} 
\]
where $\wt{\Scr C} \to S$ is a flat family of connected nodal projective curves of the
same arithmetic genus as $C$, the map $\wt{\Scr C} \to C\xt S$ has degree $1$, 
and the induced section $\wt{\Scr C} \to \sP_G \xt^G G/B$ is \emph{stable} 
in the sense of \cite{Kontsevich} over every geometric point of $S$.
Let $\mathscr M_C$ denote the moduli stack of proper connected nodal curves $\wt{\Scr C}$
equipped with a degree one map to $C$. 
Then the natural map $\ol\Bun_B^K\to \mathscr M_C$ is smooth, and $\mathscr M_C$ is
smooth (\cite[Proposition 2.4.1]{Campbell}) and there is a proper map
$\ol\Bun_B^K \to \ol\Bun_B$ over $\Bun_G$ (\cite[Proposition 3.2.2]{Campbell}). 
Inside $\mathscr M_C$ we have the open point corresponding to $\mathrm{id} : C\to C$
and the complement $\partial\mathscr M_C$ is a normal crossings divisor (\cite[Proposition 4.4.1]{Campbell}). Evidently the preimage of the open point in $\ol\Bun_B^K$ identifies
with $\Bun_B$. 
Let $_{\le\ch\nu}\ol\Bun_B^{K,-\ch\mu}$ denote the preimage of $_{\le\ch\nu}\ol\Bun_B^{-\ch\mu}$. Then \cite[Proposition 4.1.1]{Campbell} shows that
for $\ch\nu$ fixed and $\ch\mu$ large enough, the map 
$_{\le\ch\nu}\ol\Bun_B^{K,-\ch\mu} \to \mathscr M_C \xt \Bun_G$ is smooth. 
Now for any $\sP_G \in \Bun_G(k)$, the fiber of $\wt{\mf p}^K: {_{\le\ch\nu}}\ol\Bun_B^{K,-\ch\mu} \to \Bun_G$ is smooth over $\mathscr M_C$.
The preimage over the open point equals $(\wt{\mf p}^K)^{-1}(\sP_G) \cap \Bun_B$, 
and since $\partial \mathscr M_C$ is a normal crossings divisor, 
$(\wt{\mf p}^K)^{-1}(\sP_G) \cap \Bun_B$ must be dense in $(\wt{\mf p}^K)^{-1}(\sP_G)$.
Since $(\wt{\mf p}^K)^{-1}(\sP_G)$ is a resolution of
the fiber of $_{\le\ch\nu}\ol\Bun_B^{-\ch\mu}\to\Bun_G$, we have proved (ii).
\end{proof}

\begin{cor}
\source{intersection-complexes-unramified-l-factors}
\notready \label{cor:barYdense}
\source{intersection-complexes-unramified-l-factors}
For any $\ch\lambda\in \mathfrak c_X$,
the subscheme $\sY^{\ch\lambda}$ is dense in $\barY^{\ch\lambda}$. 
\end{cor}

\begin{proof}
The subschemes $_{\le \ch\eta}\barY^{\ch\lambda} = \barY^{\ch\lambda} \xt_{\ol\Bun_B} {_{\le\ch\eta}}\ol\Bun_B$
for $\ch\eta\in \ch\Lambda^\pos_G$ form an open covering of $\barY^{\ch\lambda}$. 
Since $\barY^{\ch\lambda}$ is quasicompact, there must exist some 
$\ch\eta$ such that $_{\le \ch\eta}\barY^{\ch\lambda} = \barY^{\ch\lambda}$.
Fix $\ch\lambda,\ch\eta$ as above. Then we can choose 
$\ch\mu \in \ch\Lambda^\pos_G$ such that $\ch\lambda+\ch\mu-\ch\nu$ is 
large enough, for all $0 \le \ch\nu\le \ch\eta$, for the purposes of
Proposition~\ref{prop:ULA} and Lemma~\ref{lem:BunBtoGsm}. 
Now we may consider 
$_{\le\ch\eta}\barY^{\ch\lambda+\ch\mu}$ as an open subscheme of 
$Y':=\Scr M_X \underset{\Bun_G}\xt {_{\le \ch\eta}\ol\Bun^{-\ch\lambda-\ch\mu}_B}$.
Proposition~\ref{prop:ULA}(ii) implies that 
$\sM_X \underset{\Bun_G}\xt \Bun_B^{-\ch\lambda-\ch\mu}$
is dense in $Y'$,
which implies that $\sY^{\ch\lambda+\ch\mu}$ is dense in $\barY^{\ch\lambda+\ch\mu}$.
The graded factorization property of $\barY$ 
gives a natural \'etale map 
\begin{equation} \label{e:barYfactorlarge}
\source{intersection-complexes-unramified-l-factors}
 \barY^{\ch\lambda} \oo\xt \Scr Y^{\ch\mu} := (\barY^{\ch\lambda} \xt \Scr Y^{\ch\mu})|_{\Scr A^{\ch\lambda} \oo\xt \Scr A^{\ch\mu}} \to 
\barY^{\ch\lambda+\ch\mu}, 
\end{equation}
where $\Scr Y^{\ch\mu} \into \barY^{\ch\mu}$
is the open embedding. We deduce that $\sY^{\ch\lambda} \oo\xt \sY^{\ch\mu}$ is
dense in $\barY^{\ch\lambda} \oo\xt \sY^{\ch\mu}$, which implies 
$\sY^{\ch\lambda}$ is dense in $\barY^{\ch\lambda}$. 

\end{proof}

\begin{proof}[Proof of Proposition~\ref{prop:ICbarY}]

Applying Lemma~\ref{lem:ULAbox} to the Cartesian square
\[
\begin{tikzcd} 
Y' \ar[r] \ar[d] & \Scr M_X \ar[d] \\ 
{}_{\le \ch\eta}\ol\Bun_B^{-\ch\lambda-\ch\mu} \ar[r] & \Bun_G,
\end{tikzcd}
\]
we can identify
\begin{equation} \label{e:ICbar=bt}
\source{intersection-complexes-unramified-l-factors}
    \IC_{Y'} \cong \IC_{\Scr M_X} \underset{\Bun_G}\boxtimes 
\IC_{_{\le\ch\eta} \ol\Bun^{-\ch\lambda-\ch\mu}_B}. 
\end{equation}
In particular, this gives us a description of $\IC_{_{\le\ch\eta}\barY^{\ch\lambda+\ch\mu}} = \IC_{Y'}|_{_{\le\ch\eta}\barY^{\ch\lambda+\ch\mu}}$. 
For $0 \le \ch\nu \le \ch\eta$, we have a Cartesian square
\[ 
\begin{tikzcd}
C_{\ch\nu} \times \Scr Y^{\ch\lambda+\ch\mu-\ch\nu} \ar[r,hook,"\mf i_{\barY,\ch\nu}"] \ar[d] & 
\barY^{\ch\lambda+\ch\mu} \ar[d] \\ 
C_{\ch\nu} \times \Bun_B^{-\ch\lambda-\ch\mu+\ch\nu} \ar[r,hook,"\mf i_{\ch\nu}"] &
\ol\Bun_B^{-\ch\lambda-\ch\mu} 
\end{tikzcd}
\]

Theorem \ref{thm:BFGM}, together with the Cartesian square above and 
\eqref{e:ICbar=bt}, allow us to deduce that there exists an isomorphism 
\[ \mf i_{\barY,\ch\nu}^*( \IC_{\barY^{\ch\lambda+\ch\mu}} ) 
\cong \mfU^\vee(\ch{\mfn}_C)^{-\ch\nu} \bt (\IC_{\Scr M_X} \bt_{\Bun_G} \IC_{\Bun_B^{-\ch\lambda-\ch\mu+\ch\nu}})|_{\Scr Y^{\ch\lambda+\ch\mu-\ch\nu}}
= \mfU^\vee(\ch{\mfn}_C)^{-\ch\nu} \bt (\IC_{\Scr M_X}|^{!*}_{\Scr Y^{\ch\lambda+\ch\mu-\ch\nu}}) \]
on $C_{\ch\nu} \xt \Scr Y^{\ch\lambda+\ch\mu-\ch\nu}$. 
In the last equality we have used the fact that $\Bun_B$ is smooth. 
Since we chose $\ch\lambda+\ch\mu-\ch\nu$ large enough to satisfy
Lemma~\ref{lem:BunBtoGsm}, the map $\Scr Y^{\ch\lambda+\ch\mu-\ch\nu} \to \Scr M_X$ is smooth. Therefore, $\IC_{\Scr M_X}|^{!*}_{\Scr Y^{\ch\lambda+\ch\mu-\ch\nu}} \cong \IC_{\Scr Y^{\ch\lambda+\ch\mu-\ch\nu}}$ and 
we get a canonical isomorphism 
\begin{equation} \label{e:ICbarYrestrict}
\source{intersection-complexes-unramified-l-factors}
    \mf i_{\barY,\ch\nu}^*(\IC_{\barY^{\ch\lambda+\ch\mu}}) 
    \cong \mfU^\vee(\ch\mfn_C)^{\ch\nu} \bt \IC_{\Scr Y^{\ch\lambda+\ch\mu-\ch\nu}}. 
\end{equation}

Observe that the following diagram is Cartesian:
\[
\begin{tikzcd}[column sep=large]
    (C_{\ch\nu} \xt \Scr Y^{\ch\lambda-\ch\nu}) \oo\xt \Scr Y^{\ch\mu} 
\ar[r,hook,"\mf i_{\barY,\ch\nu} \times \on{id}"] \ar[d] & 
    \barY^{\ch\lambda} \oo\xt \Scr Y^{\ch\mu} \ar[d, "\text{\eqref{e:barYfactorlarge}}"] \\ 
    C_{\ch\nu} \xt \Scr Y^{\ch\lambda+\ch\mu-\ch\nu} \ar[r, hook, "\mf i_{\barY,\ch\nu}"] & 
    \barY^{\ch\lambda+\ch\mu} 
\end{tikzcd}
\]
where the left vertical arrow is identity on $C_{\ch\nu}$ times 
the map $\Scr Y^{\ch\lambda-\ch\nu} \oo\xt \Scr Y^{\ch\mu} \to \Scr Y^{\ch\lambda+\ch\mu-\ch\nu}$  coming from graded factorization. 
Since the restriction of $\IC_{\barY^{\ch\lambda+\ch\mu}}$ to 
$\barY^{\ch\lambda} \oo\xt \Scr Y^{\ch\mu}$ is 
$\IC_{\barY^{\ch\lambda}} \mathbin{\oo\bt} \IC_{\Scr Y^{\ch\mu}}$, 
we deduce from the Cartesian square and \eqref{e:ICbarYrestrict} 
that there is a canonical isomorphism 
\[ \mf i^*_{\barY,\ch\nu}(\IC_{\barY^{\ch\lambda}}) \mathbin{\oo\bt} \IC_{\Scr Y^{\ch\mu}}
\cong (\mfU^\vee(\ch\mfn_C)^{-\ch\nu} \bt \IC_{\Scr Y^{\ch\lambda-\ch\nu}}) 
\mathbin{\oo\bt} \IC_{\Scr Y^{\ch\mu}} 
\]
when restricted to $(C_{\ch\nu} \xt \Scr Y^{\ch\lambda-\ch\nu})\oo\xt \Scr Y^{\ch\mu}$.
Lastly, let $\mathscr Y$ be a connected component of $\sY^{\ch\mu,0}$ so 
that the projection $p : (C_{\ch\nu} \xt \sY^{\ch\lambda-\ch\nu}) \oo\xt \mathscr Y \to 
C_{\ch\nu} \xt \sY^{\ch\lambda-\ch\nu}$ 
is smooth surjective with irreducible fibers. 
We have constructed an isomorphism 
\[ p^*( \mf i^*_{\barY,\ch\nu}(\IC_{\barY^{\ch\lambda}}) )
\cong p^*(  \mfU^\vee(\ch\mfn_C)^{-\ch\nu} \bt \IC_{\Scr Y^{\ch\lambda-\ch\nu}} ), \]
which implies Proposition~\ref{prop:ICbarY} because 
the functor $p^*[\dim \mathscr Y]$ is fully faithful on the category
of perverse sheaves (\cite[Proposition 4.2.5]{BBDG}).
% $\mfU^\vee(\ch\mfn_C)^{-\ch\nu}$ is NOT perverse
\end{proof}


\subsubsection{Convolution product}  \label{sect:convolution-product}
\source{intersection-complexes-unramified-l-factors}

The toric variety $X/\!\!/N$ has the natural structure of a commutative 
algebraic monoid. The multiplication operator on $X/\!\!/N$ induces
a finite map 
\[ m_{\Scr A} : \Scr A^{\ch\lambda_1} \xt \Scr A^{\ch\lambda_2} \to \Scr A^{\ch\lambda_1+\ch\lambda_2}. \]
If we have sheaves $\Scr F_i \in \rmD^b_c(\Scr A^{\ch\lambda_i}),\, i=1,2$,
we define their convolution by 
\[ \Scr F_1 \star \Scr F_2 := m_{\Scr A,!}(\Scr F_1 \bt \Scr F_2) \in 
\rmD^b_c(\Scr A^{\ch\lambda_1+\ch\lambda_2}). \]



\subsubsection{}
We have a closed embedding $\mf i_{\Scr A,\ch\nu} : C_{\ch\nu} \into \Scr A^{\ch\nu}$ corresponding to the partition $\sum_i n_i [\ch\alpha_i]$
of degree $\ch\nu = \sum_i n_i\ch\alpha_i$. 


\begin{cor}
\source{intersection-complexes-unramified-l-factors}
\notready \label{cor:pibarY}
\source{intersection-complexes-unramified-l-factors}
There is an equality 
\[ [ \pi_!(\IC_{\Scr Y^{\ch\lambda}}) ] = \sum_{\ch\nu\in \ch\Lambda^\pos_G} 
[ \mf i_{\Scr A,\ch\nu,!}(\Upsilon(\ch\mfn_C)^{\ch\nu}) \star \bar\pi_!(\IC_{\barY^{\ch\lambda-\ch\nu}})] \]
in the Grothendieck group of perverse sheaves on $\Scr A^{\ch\lambda}$. 
\end{cor}

\begin{proof}
Taking the Grothendieck--Cousin complex associated to the stratifications
$\mf i_{\barY,\ch\nu}$ and applying Proposition~\ref{prop:ICbarY} 
gives an equality 
\[ [\IC_{\barY^{\ch\lambda}}] = \sum_{\ch\nu \in \ch\Lambda^\pos_G} 
[\mf i_{\barY,\ch\nu,!} ( \mf U^\vee(\ch\mfn_C)^{\ch\nu} \bt \IC_{\Scr Y^{\ch\lambda-\ch\nu}} ) ] \]
in the Grothendieck group of perverse sheaves on $\barY^{\ch\lambda}$.
Note that $\Scr Y^{\ch\lambda-\ch\nu}$ is nonempty for finitely many
values of $\ch\nu$. 
The composition $\bar\pi \circ \mf i_{\barY,\ch\nu} : C_{\ch\nu} \xt \Scr Y^{\ch\lambda-\ch\nu} \to \Scr A^{\ch\lambda}$ coincides with
the composition $C_{\ch\nu} \xt \Scr Y^{\ch\lambda-\ch\nu} \overset{\mf i_{\Scr A,\ch\nu} \xt \pi}\longrightarrow \Scr A^{\ch\nu} \xt \Scr A^{\ch\lambda-\ch\nu} \overset{m_{\Scr A}}\to \Scr A^{\ch\lambda}$.
Therefore, applying $\bar\pi_!$ to the equality above, we get
the equality 
\[ [\bar\pi_!(\IC_{\barY^{\ch\lambda}})] = \sum_{\ch\nu\in \ch\Lambda^\pos_G} [\mf i_{\Scr A,\ch\nu,!}(\mf U^\vee(\ch\mfn_C)^{\ch\nu}) \star \pi_!(\IC_{\Scr Y^{\ch\lambda-\ch\nu}})]  \]
in the Grothendieck group of perverse sheaves on $\Scr A^{\ch\lambda}$.
Applying a further convolution by any $\Scr T \in \rmD^b_c(C_{\ch\nu'}), 
\ch\nu'\in \ch\Lambda^\pos_G$
gives 
\[ [ \mf i_{\Scr A, \ch\nu',!}( \Scr T ) \star \bar\pi_! (\IC_{\barY^{\ch\lambda}}) ] = \sum_{\ch\nu} [ \mf i_{\Scr A,\ch\nu+\ch\nu',!}( \Scr T \star \mf U^\vee(\ch\mfn_C)^{\ch\nu}) \star \pi_!(\IC_{\Scr Y^{\ch\lambda-\ch\nu}}) ] . \]
It is known (\cite[\S 6.4]{BG2}) that for a fixed nonzero $\ch\nu \in \ch\Lambda^\pos_G$, we have an equality 
\[ \sum_{\substack{\ch\nu_1,\ch\nu_2 \in \ch\Lambda^\pos_G\\ 
\ch\nu_1+\ch\nu_2=\ch\nu}} 
[ \Upsilon(\ch\mfn_C)^{\ch\nu_1} \star \mf U^\vee(\ch\mfn_C)^{-\ch\nu_2}  ] = 0 \]
in the Grothendieck group of perverse sheaves on $C_{\ch\nu}$. 
The two preceeding equalities and induction prove the claim.
\end{proof}
